\begin{abstract}

El diseño del calendario de horarios de clases a nivel universitario (University Course Timetabling Problem) es un problema de optimización combinatoria, por lo cual escala notoriamente en complejidad junto a las dimensiones de la institución (cantidad de cursos, docentes, estudiantes, programas de estudio). En este proyecto se aborda la asignación de horarios considerando las preferencias de los docentes para una grilla de bloques horarios predefinidos, con un horizonte semanal. Este trabajo es la continuación de lo realizado en los últimos dos años para resolver este problema en la Facultad de Ingeniería de la Universidad de Montevideo (FIUM), donde se ha utilizado un método exacto basado en el software Gurobi para resolver problemas de programación entera-mixta (MIP). Los métodos exactos permiten encontrar soluciones óptimas, aunque pueden ser computacionalmente costosos para instancias grandes. Por otro lado, los métodos heurísticos, como algoritmos genéticos, búsqueda tabú y recocido simulado, ofrecen soluciones eficientes y cercanas al óptimo en tiempos razonables, siendo más adecuados para problemas de gran escala. Con esta alternativa, se pretende mejorar los tiempos de obtención de soluciones y la escalabilidad. Se presenta una comparación de los enfoques exacto y heurístico, además de un análisis multiobjetivo en el que se evalúan condiciones entre la superposición de cursos, utilización de horarios excepcionales fuera del rango habitual y la cantidad de días de asistencia semanales de los docentes.

\end{abstract}

